
% =====================================================================
\section{The ground premise: identifiability, the whole, and the derived floor}
\label{sec:ground}
% =====================================================================

The Finiteness Premise (\Cref{ax:finite}) posits a positive floor. We now show
the floor need not be posited: it is forced by a more primitive fact---that a
thing can be identified at all within a whole that is not completable. This
section derives the positive floor from identifiability and non-completability,
so that \Cref{ax:finite} becomes a consequence rather than an assumption, and
identifies the conserved reference---the \emph{medium}---against which all
individuation occurs.

\subsection{The whole and its parts}

\begin{definition}[The whole; a part]
\label{def:whole}
The \emph{whole} is the total system $\med$ within which individuation occurs.
A \emph{thing} is a part $A$ of $\med$ that is \emph{proper}: $A\neq\med$ and
$A\neq\varnothing$. A thing is thus, by definition, not the whole---anything that
were the whole would not be a distinct thing within it but the whole itself.
\end{definition}

\begin{definition}[Non-completable whole]
\label{def:noncompletable}
The whole $\med$ is \emph{non-completable} if no finite stage exhausts it: for
every finite sub-collection $F$ of parts already individuated, there is a further
part of $\med$ not in $F$. We do not assume $\med$ is a completed infinite set;
we assume only that the process of exhausting it does not terminate. This is the
sole sense of ``infinite'' used in the paper, and it is an order condition
\textup{(}no terminal stage\textup{)}, not a cardinality posit.
\end{definition}

\begin{remark}[Why non-completability, not cardinality]
\label{rem:noncompletable}
We are careful to phrase the inexhaustibility of the whole as
non-completability rather than as a completed uncountable set, for two reasons.
First, the substructures of this paper \textup{(}attained minimum cuts, a
positive floor, determinate complements\textup{)} require finiteness at each
stage and could fail against a completed infinity \textup{(}\Cref{sec:master}\textup{)}.
Second, and decisively, it is exactly non-completability---the impossibility of
finishing the comparison of a part against the whole---that forces the positive
floor below. A completed totality is not needed and would be the wrong object; a
never-finishing process is both sufficient and correct.
\end{remark}

\subsection{Identification and its residual}

\begin{definition}[Identification]
\label{def:identification}
To \emph{identify} a thing $A$ is to compare it against the rest of the whole:
to determine $A$ by individuating it from $\med\setminus A$
\textup{(}individuation by negation, \Cref{thm:negation}\textup{)}. An
identification is \emph{complete} if the comparison of $A$ against the entirety of
$\med\setminus A$ is finished---every part of the complement has been brought to
bear in distinguishing $A$.
\end{definition}

\begin{definition}[Identification residual]
\label{def:idresidual}
The \emph{identification residual} of $A$, written $\idres(A)$, is the content of
the comparison of $A$ against $\med\setminus A$ that remains unfinished. A
complete identification has $\idres(A)=0$; an identification that cannot be
finished has $\idres(A)>0$, the irreducible remainder.
\end{definition}

\subsection{The floor is forced by infinitude}

\begin{theorem}[Floor from Infinitude]
\label{thm:floor-infinitude}
Let $\med$ be a non-completable whole \textup{(}\Cref{def:noncompletable}\textup{)}
and let $A$ be a thing \textup{(}a proper part, \Cref{def:whole}\textup{)} that is
identified by comparison against $\med\setminus A$. Then the identification of $A$
cannot be completed, and its identification residual is strictly positive:
\[
\boxed{\ \idres(A)\ >\ 0.\ }
\]
\end{theorem}

\begin{proof}
Suppose, for contradiction, that the identification of $A$ were complete, so that
$\idres(A)=0$. By \Cref{def:identification}, completion means the comparison of
$A$ against the entirety of $\med\setminus A$ is finished: every part of the
complement has been brought to bear. But to finish comparing $A$ against
$\med\setminus A$ is to exhaust $\med\setminus A$, and since $A$ is a fixed proper
part, exhausting $\med\setminus A$ exhausts $\med$ up to the fixed $A$---a finite
addition. Hence a completed identification yields a finite stage that exhausts
$\med$, contradicting non-completability \textup{(}\Cref{def:noncompletable}\textup{)}.
Therefore the identification cannot be completed, and $\idres(A)>0$.
\end{proof}

\begin{corollary}[The positive floor is derived]
\label{cor:floor-derived}
For a thing $A$ in a non-completable whole, the irreducible boundary content of
its individuation is bounded below by a positive constant
$\floorc:=\inf_A\idres(A)>0$. The positive floor of \Cref{ax:finite} is therefore
not an independent assumption but a consequence of: \textup{(}i\textup{)} a thing
is a proper part, and \textup{(}ii\textup{)} the whole is non-completable.
\end{corollary}

\begin{proof}
By \Cref{thm:floor-infinitude}, $\idres(A)>0$ for every thing $A$; the infimum of
strictly positive identification residuals over the things of $\med$ is a
constant $\floorc\ge0$. If $\floorc=0$ there would be things with arbitrarily
small residual, i.e. arbitrarily near-complete identifications; but each such
identification, in the limit, would finish the comparison and exhaust $\med$, the
same contradiction as in \Cref{thm:floor-infinitude}. Hence $\floorc>0$. Taking
this $\floorc$ as the edge-weight floor recovers \Cref{ax:finite} as a derived
statement: every separating boundary in the contact graph of $\med$'s parts costs
at least the unfinishable remainder $\floorc>0$.
\end{proof}

\begin{remark}[The floor is the trace of the whole on its parts]
\label{rem:floor-trace}
\Cref{thm:floor-infinitude} reverses the status of the floor. It is not a
limitation of measurement nor a stipulated constant; it is the structural
shadow that a non-completable whole casts on each of its finite parts. A finite
thing in an inexhaustible whole cannot complete its comparison against the rest,
and the unfinishable remainder of that comparison \emph{is} the floor. Were the
whole completable \textup{(}finite\textup{)}, the comparison could finish, the
residual could reach zero, and a thing's truth could be a point; because the
whole is not completable, the residual is positive and a thing's truth is a
region \textup{(}\Cref{sec:residual}\textup{)}. Finitude makes the thing;
non-completability of the whole makes the thing's truth a cell.
\end{remark}

\subsection{The medium as conserved reference}

Identification is comparison against the rest. We name the reference and record
that it is the conserved invariant of the construction.

\begin{definition}[Medium]
\label{def:medium}
The \emph{medium} is the whole $\med$ in its role as the reference against which
every part is individuated: the totality simultaneously in contact with every
part, such that the separation cost of a part $A$ is the cost of its boundary to
$\med\setminus A$. Individuation of $A$ is undefined absent the medium---there is
then no reference against which $A$ is a part rather than the whole.
\end{definition}

\begin{proposition}[Individuation requires the medium]
\label{prop:medium-required}
For a thing $A$, the separation cost $b(A)=\wt(\cut(A))$ is defined only relative
to the medium $\med\setminus A$. Removing the medium leaves $b(A)$ undefined:
without a complement to be individuated against, $A$ is not determined as a part.
\end{proposition}

\begin{proof}
By \Cref{thm:negation} a part is fixed only as the complement of its
negation-set; the negation-set of $A$ within $\med$ is $\med\setminus A$, the
medium relative to $A$. The boundary $\cut(A)$ is the set of edges between $A$ and
$\med\setminus A$; with no $\med\setminus A$ there are no such edges and $b(A)$
has no value. Hence individuation is defined only against the medium.
\end{proof}

\begin{remark}[Marbles in water]
\label{rem:marbles}
The relation of parts to medium is that of marbles to the water they sit in. The
marbles are the individuated parts---what is known, the knowledge built so far;
the water is the medium---the conserved invariant in contact with every marble.
Rearranging the marbles \textup{(}building new boundaries, re-encoding\textup{)}
does not change the water: the medium is invariant under all such reshuffling
\textup{(}\Cref{thm:residual-invariant}\textup{)}. A thing's identity is the
water on its own side of its boundary---the conserved residue it holds---not the
marbles around it nor the boundary edges that mark it off. We return to this
reading, on which no theorem depends, in \Cref{sec:measurement}.
\end{remark}

\begin{remark}[Identifiability is the premise of the whole ladder]
\label{rem:identifiability}
\Cref{thm:floor-infinitude} and \Cref{cor:floor-derived} let us read the entire
paper from a single near-undeniable starting point: \emph{that anything is
identifiable at all within a non-completable whole}. From this the positive floor
\textup{(}T0\textup{)} follows by \Cref{cor:floor-derived}; individuation by
negation \textup{(}T1\textup{)} follows because identification is comparison
against the rest \textup{(}\Cref{def:identification}\textup{)}; and the remaining
rungs T2--T8 follow from these as proved below. The Finiteness Premise is retained
as the working hypothesis for the graph-theoretic development, now justified as a
consequence of identifiability rather than an axiom in its own right. We make this
equivalence precise in the Master Theorem \textup{(}\Cref{sec:master}\textup{)}.
\end{remark}
